% Worked solutions explicitly requested by the user.
\documentclass[11pt]{article}
\usepackage{../../eecs245}
\enablesolutions
\enablestudentinfo
\enableexam
\usepackage{amsmath,amsfonts,amssymb,multicol}
\DeclareSymbolFont{operators}{OT1}{cmr}{m}{n}
\DeclareSymbolFont{letters}{OML}{cmm}{m}{it}
\DeclareSymbolFont{symbols}{OMS}{cmsy}{m}{n}
\raggedbottom
\newcommand{\coursename}{EECS 245}
\newcommand{\term}{Fall 2026 at the University of Michigan}
\newcommand{\assignment}{Midterm 2}
\newcommand{\submissioninstructions}{%
\textbf{Instructions}
\begin{itemize}
\item The student exam consists of 7 problems, worth a total of 100 points, spread across 14 pages (7 sheets of paper).
\item You have 120 minutes to complete this exam, unless you have extended-time accommodations through SSD.
\item Write your uniqname in the top right corner of each page.
\item For free response problems, \textbf{you must show all of your work}, and write your final answer in the space provided. We will not grade work that appears elsewhere, and you may lose points if your work is not shown.
\item For multiple choice problems, completely fill in bubbles and square boxes; if we cannot tell which option(s) you selected, you may lose points.\\
\bubble{A bubble means that you should only select one choice.}\\
\squarebubble{A square box means you should select all that apply.}
\item You may refer to \textbf{two double-sided 8.5x11" handwritten notes sheets}. Other than that, you may not refer to any other resources or technology during the exam (no phones, watches, or calculators).
\end{itemize}}
\newcommand{\pairanswerbox}[4][1cm]{%
  \noindent\begin{tikzpicture}[baseline=(ansbox.south)]
    \node[draw, inner sep=0pt, minimum width=\dimexpr\linewidth-0.8pt\relax, minimum height=#4, anchor=north west] (ansbox) at (0,0) {};
    \node[anchor=south east, xshift=-6pt, yshift=6pt] at (ansbox.south east) {$#2=\minibox{2.3cm}{}[#1]\qquad #3=\minibox{2.3cm}{}[#1]$};
  \end{tikzpicture}%
}
\newcommand{\setanswerbox}[4]{%
  \noindent\begin{tikzpicture}[baseline=(ansbox.south)]
    \node[draw, inner sep=0pt, minimum width=\dimexpr\linewidth-0.8pt\relax, minimum height=#2, anchor=north west] (ansbox) at (0,0) {};
    \node[anchor=south east, xshift=-6pt, yshift=6pt] at (ansbox.south east) {$#1=\left\{\minibox{#3}{}[#4]\right\}$};
  \end{tikzpicture}%
}
\newcommand{\workbox}[1]{\par\noindent\emptybox[\dimexpr\linewidth-2\fboxsep-2\fboxrule\relax]{#1}\par}
\newcommand{\projectionoptionstrut}{\rule[-0.45cm]{0pt}{1.15cm}}
\begin{document}
\makemytitle{\coursename}{\term}{\assignment}{}{\submissioninstructions}{\bubble{1013 DOW} \bubble{1014 DOW} \bubble{4721 BBB} \bubble{Other}}
\vspace{1em}
You are to abide by the University of Michigan/Engineering Honor Code. To receive a grade, please sign below to signify that you have kept the Honor Code pledge.

\textit{I have neither given nor received aid on this exam, nor have I concealed any violations of the Honor Code.}

\emptybox{1.5cm}
\newpage
\begin{center}
\null\vspace{9cm}
This page has been intentionally left blank.
\end{center}
\newpage

\textbf{Tip: Skim through the entire exam before starting to work on it.}

\begin{prob}[(10 pts)]
Let
\[
A=\begin{bmatrix}6&\alpha\\8&12\end{bmatrix}
\]
Each part below is independent. Find the value of $\alpha$ that satisfies the given condition.
\begin{subprob}(3 pts)
$\det(A)=40$.\par\vspace{0.2em}

\begin{solution}
For this $2\times2$ matrix,
\[
\det(A)=6(12)-8\alpha=72-8\alpha
\]
Thus,
\[
72-8\alpha=40\quad\Longrightarrow\quad\boxed{\alpha=4}
\]
\end{solution}
\end{subprob}
\begin{subprob}(3 pts)
$\text{rank}(A)=1$.\par\vspace{0.2em}

\begin{solution}
Since $A$ is a nonzero $2\times2$ matrix, it has rank $1$ exactly when its determinant is $0$. Therefore,
\[
72-8\alpha=0\quad\Longrightarrow\quad\boxed{\alpha=9}
\]
For this value, the second column is $3/2$ times the first column.
\end{solution}
\end{subprob}
\begin{subprob}(4 pts)
$A^{-1}=\begin{bmatrix}3/2&-1\\-1&3/4\end{bmatrix}$.\par\vspace{0.2em}

\begin{solution}
The inverse formula gives
\[
A^{-1}=\frac{1}{72-8\alpha}\begin{bmatrix}12&-\alpha\\-8&6\end{bmatrix}
\]
Matching the top-left entry gives
\[
\frac{12}{72-8\alpha}=\frac32
\quad\Longrightarrow\quad72-8\alpha=8
\quad\Longrightarrow\quad\boxed{\alpha=8}
\]
Substituting $\alpha=8$ gives all four entries of the specified inverse.
\end{solution}
\end{subprob}
\end{prob}
\newpage

\begin{prob}[(10 pts)]
Let
\[
B=\begin{bmatrix}2&2&2\\3&-3&0\\1&1&-2\end{bmatrix}
\]
\begin{subprob}(4 pts)
Find $BB^T$. \\

\begin{solution}
The entries of $BB^T$ are the dot products of the rows of $B$. The three rows are mutually orthogonal, and their squared lengths are
\[
2^2+2^2+2^2=12,\qquad 3^2+(-3)^2=18,\qquad 1^2+1^2+(-2)^2=6
\]
Hence,
\[
\boxed{BB^T=\begin{bmatrix}12&0&0\\0&18&0\\0&0&6\end{bmatrix}}
\]
\end{solution}
\end{subprob}
\begin{subprob}(6 pts)
Find $B^{-1}$. Fully simplify each entry. \textit{Hint: Use part (a), not $\det(B)$ or a $3\times3$ inverse formula. If $BB^T=C$, what is $B^{-1}$?}\par\vspace{0.2em}

\begin{solution}
Let $C=BB^T$. Since $C$ is invertible,
\[
B(B^TC^{-1})=(BB^T)C^{-1}=CC^{-1}=I
\]
Thus, $B^{-1}=B^TC^{-1}$. Using part (a),
\[
B^{-1}
=\begin{bmatrix}2&3&1\\2&-3&1\\2&0&-2\end{bmatrix}
\begin{bmatrix}1/12&0&0\\0&1/18&0\\0&0&1/6\end{bmatrix}
=\boxed{\begin{bmatrix}1/6&1/6&1/6\\1/6&-1/6&1/6\\1/6&0&-1/3\end{bmatrix}}
\]
\end{solution}
\end{subprob}
\end{prob}
\newpage

\begin{prob}[(16 pts)]
Let $X=CR$, where $a,b$ are unknown scalars. Furthermore, suppose that
\[
C=\begin{bmatrix}4&0\\-2&5\\3&-1\\0&6\end{bmatrix},\qquad
R=\begin{bmatrix}1&3&a&-2&5\\0&4&8&b&0\end{bmatrix},\qquad
\begin{bmatrix}-5\\0\\1\\-2\\1\end{bmatrix}\in\text{nullsp}(X)
\]
\begin{subprob}(4 pts)
Determine the following values. \\

$\text{rank}(X)=$ \minibox{3cm}{$2$}[1.2cm]

$\dim(\text{nullsp}(X))=$ \minibox{3cm}{$3$}[1.2cm]

\begin{solution}
The columns of $C$ are independent, and the rows of $R$ are independent. This CR factorization therefore has $\text{rank}(X)=2$. Since $X$ has five columns, rank-nullity gives $\dim(\text{nullsp}(X))=5-2=3$.
\end{solution}
\end{subprob}
\begin{subprob}(8 pts)
Find $a$ and $b$. \\

\begin{solution}
Let $\vec v$ be the given vector in $\text{nullsp}(X)$. Then $CR\vec v=\vec0$. Since the columns of $C$ are linearly independent, this requires $R\vec v=\vec0$. Therefore,
\[
R\vec v=\begin{bmatrix}-5+a+4+5\\8-2b\end{bmatrix}
=\begin{bmatrix}a+4\\8-2b\end{bmatrix}=\vec0
\]
It follows that $\boxed{a=-4,\ b=4}$.
\end{solution}
\end{subprob}
\begin{subprob}(4 pts)
Which of the following choices of columns of $X$ form a basis for $\text{colsp}(X)$? \textbf{Select all that apply.} \\

\correctsquarebubble{Columns 1 and 2}

\squarebubble{Columns 1 and 5}

\correctsquarebubble{Columns 2 and 3}

\squarebubble{Columns 3 and 4}

\correctsquarebubble{Columns 4 and 5}

\squarebubble{Columns 1, 2, and 3}

\begin{solution}
Because $C$ has linearly independent columns, we can check independence using the corresponding columns of $R$, with $a=-4$ and $b=4$. Columns $1$ and $5$ are scalar multiples, as are columns $3$ and $4$. Each remaining listed pair is independent.

Since $\text{rank}(X)=2$, the valid bases are \textbf{columns 1 and 2, columns 2 and 3, and columns 4 and 5}. Three columns cannot form a basis for this two-dimensional space.
\end{solution}
\end{subprob}
\end{prob}
\newpage

\begin{prob}[(16 pts)]
Let
\[
A=\begin{bmatrix}
0&0&1&0\\
1&0&0&0\\
0&1&0&0\\
0&0&0&-1
\end{bmatrix},\qquad
\vec v=\begin{bmatrix}1\\2\\3\\4\end{bmatrix}
\]
\begin{subprob}(2 pts)
Find $A\vec v$.\par\vspace{0.2em}

\begin{solution}
Multiplying by $A$ cycles the first three components and negates the fourth:
\[
A\begin{bmatrix}v_1\\v_2\\v_3\\v_4\end{bmatrix}
=\begin{bmatrix}v_3\\v_1\\v_2\\-v_4\end{bmatrix}
\quad\Longrightarrow\quad
\boxed{A\vec v=\begin{bmatrix}3\\1\\2\\-4\end{bmatrix}}
\]
\end{solution}
\end{subprob}
\begin{subprob}(2 pts)
What is the smallest positive integer $k$ such that $A^k=I$?

\bubble{2}\bubble{3}\bubble{4}\correctbubble{6}\bubble{8}\bubble{12}

\begin{solution}
The first three components return to their original positions after a multiple of $3$ applications of $A$. The fourth returns to its original sign after an even number of applications. The smallest positive integer divisible by both $3$ and $2$ is $\boxed{k=6}$.
\end{solution}
\end{subprob}
\begin{subprob}(6 pts)
Give one basis for $\text{nullsp}(A-I)$.\par\vspace{0.2em}

\begin{solution}
A vector is in $\text{nullsp}(A-I)$ exactly when $A\vec x=\vec x$. Componentwise,
\[
x_3=x_1,\qquad x_1=x_2,\qquad x_2=x_3,\qquad -x_4=x_4
\]
Thus, $x_1=x_2=x_3$ and $x_4=0$. Every vector in the null space is a multiple of $\begin{bmatrix}1\\1\\1\\0\end{bmatrix}$, so one basis is
\[
\boxed{\left\{\begin{bmatrix}1\\1\\1\\0\end{bmatrix}\right\}}
\]
\end{solution}
\end{subprob}
\newpage

Recall,
\[
A=\begin{bmatrix}
0&0&1&0\\
1&0&0&0\\
0&1&0&0\\
0&0&0&-1
\end{bmatrix}
\]
\begin{subprob}(6 pts)
Applying $A$ to a vector $\vec x$ changes its components by the amounts shown below:
\[
A\vec x-\vec x=\begin{bmatrix}m\\-3\\5\\2\end{bmatrix}
\]
\begin{enumerate}[label=(\roman*)]
\item (4 pts) Find $m$. Your answer should be a number with no variables.\par\vspace{0.2em}

\vspace{0.7cm}
\begin{solution}
Applying $A$ rearranges the first three components, so their sum does not change. Their changes must therefore sum to zero:
\[
m+(-3)+5=0\quad\Longrightarrow\quad\boxed{m=-2}
\]
\end{solution}
\item (2 pts) For this value of $m$, how many vectors $\vec x$ satisfy $A\vec x-\vec x=\begin{bmatrix}m\\-3\\5\\2\end{bmatrix}$?\par\vspace{0.2em}
\bubble{0}\bubble{1}\bubble{2}\bubble{3}\bubble{4}\correctbubble{Infinitely many}

\begin{solution}
There are \textbf{infinitely many} solutions. One particular solution is $\begin{bmatrix}0\\3\\-2\\-1\end{bmatrix}$. Adding any vector in $\text{nullsp}(A-I)$ preserves the product with $A-I$, giving
\[
\vec x=\begin{bmatrix}0\\3\\-2\\-1\end{bmatrix}
+t\begin{bmatrix}1\\1\\1\\0\end{bmatrix},\qquad t\in\mathbb R
\]
\end{solution}
\end{enumerate}
\end{subprob}
\end{prob}
\newpage

\begin{prob}[(16 pts)]
Let
\[
X=\begin{bmatrix}1&4\\2&3\\2&4\\0&2\end{bmatrix},\qquad
\vec q_1=\frac13\begin{bmatrix}1\\2\\2\\0\end{bmatrix}
\]
We would like to find \textbf{orthogonal} unit vectors $\vec q_1,\vec q_2\in\mathbb R^4$ such that
{\setlength{\belowdisplayskip}{0pt}\setlength{\belowdisplayshortskip}{0pt}
\[
\text{span}(\{\vec q_1,\vec q_2\})=\text{colsp}(X)
\]}
\vspace{-0.6em}
\begin{subprob}(6 pts)
Find one possible $\vec q_2$.\par\vspace{0.2em}

\begin{solution}
Let $\vec x_2$ be the second column of $X$. Remove its component in the direction of $\vec q_1$:
\[
\vec q_1\cdot\vec x_2=\frac{4+6+8}{3}=6
\]
\[
\vec u_2=\vec x_2-6\vec q_1
=\begin{bmatrix}4\\3\\4\\2\end{bmatrix}
-\begin{bmatrix}2\\4\\4\\0\end{bmatrix}
=\begin{bmatrix}2\\-1\\0\\2\end{bmatrix}
\]
Since $\|\vec u_2\|=\sqrt{4+1+4}=3$, one answer is
\[
\boxed{\vec q_2=\frac13\begin{bmatrix}2\\-1\\0\\2\end{bmatrix}}
\]
This is a unit vector orthogonal to $\vec q_1$. The first column of $X$ is $3\vec q_1$, and its second column is $6\vec q_1+3\vec q_2$, so the two unit vectors span $\text{colsp}(X)$. The negative of this $\vec q_2$ is also valid.
\end{solution}
\end{subprob}
\newpage


For the remaining parts, suppose $Q$ has orthonormal columns $\vec q_1,\vec q_2\in\mathbb R^4$. You do not need your answer to part (a). Let
\[
Q=\begin{bmatrix}|&|\\\vec q_1&\vec q_2\\|&|\end{bmatrix},\qquad
Z=\begin{bmatrix}|&|\\\vec q_1+\vec q_2&\vec q_1-\vec q_2\\|&|\end{bmatrix}
\]
and suppose $\vec y\in\mathbb R^4$ satisfies
\[
Q^T\vec y=\begin{bmatrix}3\\1\end{bmatrix}
\]
\begin{subprob}(2 pts)
Is $\vec y\in\text{colsp}(Q)$?\par\vspace{0.2em}
\bubble{Yes}\bubble{No}\correctbubble{Impossible to tell}

\begin{solution}
\textbf{Impossible to tell.} The vector $\vec y_0=3\vec q_1+\vec q_2$ is in $\text{colsp}(Q)$ and has the specified product with $Q^T$.

There are also nonzero vectors $\vec r$ orthogonal to $\text{colsp}(Q)$, since this is a two-dimensional subspace of $\mathbb R^4$. Then $\vec y_0+\vec r$ is outside $\text{colsp}(Q)$ but has the same product with $Q^T$, because $Q^T\vec r=\vec0$. Both possibilities are consistent with the information given.
\end{solution}
\end{subprob}
\vspace{0.4cm}
\newpage
Recall that $Q$ has orthonormal columns, and
\[
Z=\begin{bmatrix}|&|\\\vec q_1+\vec q_2&\vec q_1-\vec q_2\\|&|\end{bmatrix},\qquad
Q^T\vec y=\begin{bmatrix}3\\1\end{bmatrix}
\]
\begin{subprob}(4 pts)
Which choice gives both the vector $\vec w^\ast$ that minimizes $\|\vec y-Z\vec w\|^2$ and the resulting projection, $\vec p=Z\vec w^\ast$?\par\vspace{0.2cm}

\bubble{A. $\vec w^\ast=\begin{bmatrix}3\\1\end{bmatrix},\quad \vec p=3\vec q_1+\vec q_2$}

\correctbubble{B. $\vec w^\ast=\begin{bmatrix}2\\1\end{bmatrix},\quad \vec p=3\vec q_1+\vec q_2$}

\bubble{C. $\vec w^\ast=\begin{bmatrix}3\\1\end{bmatrix},\quad \vec p=4\vec q_1+2\vec q_2$}

\bubble{D. $\vec w^\ast=\begin{bmatrix}2\\1\end{bmatrix},\quad \vec p=2\vec q_1+\vec q_2$}

\bubble{E. $\vec w^\ast=\begin{bmatrix}1\\2\end{bmatrix},\quad \vec p=3\vec q_1+\vec q_2$}

\begin{solution}
\textbf{B.} Adding and subtracting the columns of $Z$ recovers $2\vec q_1$ and $2\vec q_2$, so $Z$ and $Q$ have the same column space. Hence,
\[
\boxed{\vec p=QQ^T\vec y=3\vec q_1+\vec q_2}
\]
But $Z\vec w=(w_1+w_2)\vec q_1+(w_1-w_2)\vec q_2$. Independence of $\vec q_1,\vec q_2$ therefore requires
\[
w_1^\ast+w_2^\ast=3,\qquad w_1^\ast-w_2^\ast=1
\quad\Longrightarrow\quad
\boxed{\vec w^\ast=\begin{bmatrix}2\\1\end{bmatrix}}
\]
$Q^T\vec y$ gives coefficients in $Q$'s columns, while $\vec w^\ast$ gives coefficients in $Z$'s columns.
\end{solution}
\end{subprob}
\vspace{0.2cm}
\vspace{0.3em}
\begin{subprob}(4 pts)
Identify whether each statement is true or false. \\

\begin{enumerate}[label=(\roman*)]
\item (2 pts) For every $\vec x\in\mathbb R^2$, $Q^TQ\vec x=\vec x$.
\correctbubble{True} \bubble{False}

\begin{solution}
\textbf{True.} Orthonormal columns give $Q^TQ=I_2$, so $Q^TQ\vec x=\vec x$ for every $\vec x\in\mathbb R^2$.
\end{solution}
\item (2 pts) $QQ^T$ is invertible.
\bubble{True} \correctbubble{False}

\begin{solution}
\textbf{False.} $QQ^T$ is $4\times4$, and its columns lie in the two-dimensional space $\text{colsp}(Q)$. Its rank is at most $2$, so it is not invertible.
\end{solution}
\end{enumerate}
\end{subprob}
\end{prob}
\newpage

\begin{prob}[(14 pts)]
Suppose we'd like to fit a multiple linear regression model to predict the daily electricity use of household $i$, $y_i$, in kilowatt-hours.

For each household $i$, the feature vector is $\vec x_i=\begin{bmatrix}x_i^{(1)}\\x_i^{(2)}\end{bmatrix}$, where:
\begin{itemize}[topsep=0.2em,itemsep=0.2em]
\item $x_i^{(1)}$ is the number of residents.
\item $x_i^{(2)}$ is $2$ if the household uses a heat pump and $0$ if it uses gas heating.
\end{itemize}
Our model is
\[
h(\vec x_i)=w_0+w_1x_i^{(1)}+w_2x_i^{(2)}
\]
The dataset below has two unknown values, $a$ and $b$.
\begin{center}
\renewcommand{\arraystretch}{1.35}
\begin{tabular}{@{\hspace{0.25cm}}c@{\hspace{0.25cm}}|c|c|c}
$i$ & Residents ($x_i^{(1)}$) & Heating source ($x_i^{(2)}$) & Electricity use ($y_i$)\\\hline
1 & 2 & $a$ & 12\\
2 & 5 & 2 & 30\\
3 & 2 & 2 & 16\\
4 & 4 & 0 & $b$\\
5 & 3 & 0 & 16
\end{tabular}
\end{center}
\begin{subprob}(4 pts)
Write the design matrix $X$. Your answer should be a matrix in terms of $a$ and constants.\par\vspace{0.3cm}

\begin{solution}
The intercept contributes a column of ones, followed by the residents and heating-source features:
\[
\boxed{X=\begin{bmatrix}1&2&a\\1&5&2\\1&2&2\\1&4&0\\1&3&0\end{bmatrix}}
\]
\end{solution}
\end{subprob}
\newpage

Recall, our model is $h(\vec x_i)=w_0+w_1x_i^{(1)}+w_2x_i^{(2)}$, where $x_i^{(1)}$ is the number of residents and $x_i^{(2)}$ is $2$ for a heat pump and $0$ for gas heating.
\begin{center}
\renewcommand{\arraystretch}{1.35}
\begin{tabular}{@{\hspace{0.25cm}}c@{\hspace{0.25cm}}|c|c|c}
$i$ & Residents ($x_i^{(1)}$) & Heating source ($x_i^{(2)}$) & Electricity use ($y_i$)\\\hline
1 & 2 & $a$ & 12\\
2 & 5 & 2 & 30\\
3 & 2 & 2 & 16\\
4 & 4 & 0 & $b$\\
5 & 3 & 0 & 16
\end{tabular}
\end{center}
\begin{subprob}(10 pts)
The parameter vector that minimizes mean squared error is $\vec w^\ast=\begin{bmatrix}3\\4\\3\end{bmatrix}$.

Find $a$ and $b$. Your answers should be numbers with no variables.

\textit{Hint: What properties must the error vector satisfy?}\par\vspace{0.2em}

\begin{solution}
The predictions and errors are
\[
X\vec w^\ast=\begin{bmatrix}11+3a\\29\\17\\19\\15\end{bmatrix},\qquad
\vec e=\vec y-X\vec w^\ast=\begin{bmatrix}1-3a\\1\\-1\\b-19\\1\end{bmatrix}
\]
At a least-squares solution, $\vec e$ is orthogonal to every column of $X$. Orthogonality to the intercept column gives
\[
(1-3a)+1-1+(b-19)+1=0\quad\Longrightarrow\quad b-17-3a=0
\]
Orthogonality to the residents column gives
\[
2(1-3a)+5-2+4(b-19)+3=0\quad\Longrightarrow\quad4b-68-6a=0
\]
Substituting $b=17+3a$ into the second equation gives $6a=0$. Thus,
\[
\boxed{a=0,\qquad b=17}
\]
These values give $\vec e=\begin{bmatrix}1\\1\\-1\\-2\\1\end{bmatrix}$, which is also orthogonal to the heating-source column: $2(1)+2(-1)=0$. All the normal equations are satisfied.
\end{solution}
\end{subprob}
\end{prob}
\newpage

\begin{prob}[(18 pts)]
Let $\vec x\in\mathbb R^2$, and consider
\[
A=\begin{bmatrix}2&1\\1&2\end{bmatrix},\qquad
f(\vec x)=\left(\vec x^TA\vec x-1\right)^2
\]
\begin{subprob}(6 pts)
\begin{enumerate}[label=(\roman*)]
\item (2 pts) Which of the following is the fully simplified expression for $\nabla f(\vec x)$?\par\vspace{0.2em}
\noindent\bubble{$2(\vec x^TA\vec x-1)A\vec x$}\unskip\hfill
\bubble{$4(\vec x^TA\vec x-1)\vec x$}\unskip\hfill
\correctbubble{$4(\vec x^TA\vec x-1)A\vec x$}\unskip\hfill
\bubble{$4(\vec x^TA\vec x)A\vec x$}\unskip
\par\vspace{0.3cm}
\begin{solution}
The \textbf{third expression} is correct.
\end{solution}
\item (4 pts) Justify your answer.\par\vspace{0.2em}

\begin{solution}
Let $g(\vec x)=\vec x^TA\vec x-1$, so $f(\vec x)=g(\vec x)^2$. Since $A$ is symmetric,
\[
\nabla g(\vec x)=(A+A^T)\vec x=2A\vec x
\]
The chain rule gives
\[
\boxed{\nabla f(\vec x)=2g(\vec x)\nabla g(\vec x)=4(\vec x^TA\vec x-1)A\vec x}
\]
\end{solution}
\end{enumerate}
\end{subprob}
\vspace{0.3cm}
\begin{subprob}(4 pts)
Identify whether each statement is true or false. \textit{Hint: There is no need to think about second partial derivatives.}\par\vspace{0.3cm}
\begin{enumerate}[label=(\roman*)]
\item (2 pts) $f$ is convex.
\bubble{True} \correctbubble{False}

\begin{solution}
\textbf{False.} Consider
\[
\vec u=\begin{bmatrix}1/\sqrt2\\0\end{bmatrix},\qquad
\vec v=\begin{bmatrix}-1/\sqrt2\\0\end{bmatrix}
\]
Both satisfy $\vec x^TA\vec x=1$, so $f(\vec u)=f(\vec v)=0$.

Their midpoint is $\vec0$, but
\[
f\left(\frac{\vec u+\vec v}{2}\right)=f(\vec0)=1>
\frac{f(\vec u)+f(\vec v)}{2}=0
\]
This violates the defining inequality for a convex function.
\end{solution}
\item (2 pts) $f$ attains its global minimum at exactly one vector.
\bubble{True} \correctbubble{False}

\begin{solution}
\textbf{False.} Since $f$ is a square, it is always nonnegative. The two vectors from part (i) both attain $0$, so the global minimizer is not unique.
\end{solution}
\end{enumerate}
\end{subprob}\newpage

Recall, $\vec x\in\mathbb R^2$ and
\[
A=\begin{bmatrix}2&1\\1&2\end{bmatrix},\qquad
f(\vec x)=\left(\vec x^TA\vec x-1\right)^2
\]
\begin{subprob}(8 pts)
We run gradient descent on $f$ twice, starting from the \textbf{same} unknown vector $\vec x^{(0)}$.
\begin{itemize}[itemsep=0.6em]
\item With step size $\alpha=\frac14$, the first iteration gives $\vec x^{(1)}=\frac14\begin{bmatrix}-14\\1\end{bmatrix}$.
\item With step size $\alpha=\frac34$, the first iteration gives $\vec x^{(1)}=\frac14\begin{bmatrix}-54\\11\end{bmatrix}$.
\end{itemize}
Find $\vec x^{(0)}$. Your answer should be a vector with no variables. You can answer this part without using your answer to part (a).\par\vspace{0.2em}

\begin{solution}
Let $\vec g=\nabla f(\vec x^{(0)})$. Both runs have the same gradient at their starting point. The update rule gives
\[
\vec x^{(0)}-\frac14\vec g=\frac14\begin{bmatrix}-14\\1\end{bmatrix},\qquad
\vec x^{(0)}-\frac34\vec g=\frac14\begin{bmatrix}-54\\11\end{bmatrix}
\]
Multiply the first equation by $3$ and subtract the second to eliminate $\vec g$:
\[
2\vec x^{(0)}=\frac14\left(3\begin{bmatrix}-14\\1\end{bmatrix}-\begin{bmatrix}-54\\11\end{bmatrix}\right)
=\begin{bmatrix}3\\-2\end{bmatrix}
\]
Therefore,
\[
\boxed{\vec x^{(0)}=\begin{bmatrix}3/2\\-1\end{bmatrix}}
\]
\end{solution}
\end{subprob}
\end{prob}
\end{document}
